\documentclass[10pt,a4paper]{article}
% Compile from this folder:  tectonic main.tex   (or: pdflatex main && pdflatex main)
% (Overleaf: upload this file plus the folder results/figures/ from `python run_all.py`.)
\usepackage[margin=1.8cm]{geometry}
\usepackage{amsmath,amssymb,booktabs,graphicx,caption,subcaption,float,hyperref,xcolor}
\usepackage[T1]{fontenc}
\usepackage{lmodern}
\graphicspath{{../results/figures/}{results/figures/}}
\captionsetup{font=small,labelfont=bf}
\setlength{\parskip}{2pt}
\newcommand{\code}[1]{\texttt{\small #1}}
\definecolor{acc}{HTML}{00590D}
\hypersetup{colorlinks,linkcolor=acc,citecolor=acc,urlcolor=acc}

\title{\vspace{-1.2cm}\textbf{European Option Pricing: Binomial Trees vs Black--Scholes}\\[2pt]
\large with real NSE option data and a machine-learning choice of $\sigma$}
\author{\small
\begin{tabular}{@{}l@{\quad}l@{\qquad}l@{\quad}l@{}}
Adwait Eklavya & 2024ADPS0825G & Shivank Mishra & 2024ADPS0823G\\
Mohammad Zaid Agha & 2024ADPS0840G & Pradyumn Venkata Rayasam & 2023B5AD0973G\\
Asvin Sunderiyal & 2024ADPS0867G & &
\end{tabular}\\[4pt]
\small MAC F244 \textperiodcentered{} Final report \textperiodcentered{} 15 November 2026}
\date{}

\begin{document}
\maketitle
\vspace{-1.2cm}

\section{Introduction}
European options can be priced by a discrete no-arbitrage argument (the binomial tree) or by continuous-time
stochastic calculus (Black--Scholes, BS)~\cite{bs,crr,hull}. We ask: \emph{for European options on Indian underlyings, how do the two
models compare in accuracy, speed, flexibility and real-market fit, and does the choice of model or the choice of
inputs (especially volatility) matter more?} Since the tree is a discretisation of the same GBM model, we answer
the second part by separating three error sources: numerical (tree vs BS), input (which $\sigma$) and model
(what one constant $\sigma$ cannot fit). Beyond the proposal, we extended the study to real NSE option prices,
implied volatility, and a machine-learning (ML) model that chooses the volatility input.
\textbf{Answer} (Table~\ref{tab:rule}): use Black--Scholes for European options, and a Leisen--Reimer tree only for
a cash dividend modelled from spot, early exercise or path dependence; the two differ by 0.05\,bp on NSE data, while
the choice of $\sigma$ moves prices by 11\,bp, so the $\sigma$ (implied volatility where available) matters most. All code is in Python; every number and figure below
is regenerated by \code{python run\_all.py}, and each equation is implemented by the function named next to it.

\section{Models}
\textbf{Binomial tree} (\code{src/core/binomial.py}). Split $[0,T]$ into $N$ steps, $\Delta t=T/N$. A one-period
portfolio of $\Delta=(V_u-V_d)/(S u-S d)$ shares and $B$ in the bond replicates the option in both states, so its
cost $\Delta S+B$ is the price; rewriting it gives risk-neutral valuation
$V_{n,j}=e^{-r\Delta t}[qV_{n+1,j+1}+(1-q)V_{n+1,j}]$ with $q=(e^{(r-\delta)\Delta t}-d)/(u-d)$, which lies in
$(0,1)$ iff $d<e^{(r-\delta)\Delta t}<u$ (otherwise the stock or the bond dominates: an arbitrage, shown for
$r=50\%$, $\sigma=5\%$, $N=1$ where $q=6.97$). CRR~\cite{crr} uses $u=1/d=e^{\sigma\sqrt{\Delta t}}$; Jarrow--Rudd (JR)~\cite{jr} centres
the log-steps on the drift; Leisen--Reimer (LR, odd $N$)~\cite{lr} chooses $q=h(d_2)$, $q'=h(d_1)$ by Peizer--Pratt inversion
so the tree matches the BS probabilities. Tree Greeks come from the nodes at steps 1--2 (\code{binomial\_greeks}).

\textbf{Black--Scholes} (\code{src/core/bs.py}). Under $\mathbb Q$, $dS_t=(r-\delta)S_t\,dt+\sigma S_t\,dW_t$. Itô's
lemma on $V(S,t)$ and a delta-hedged portfolio give
$V_t+\tfrac12\sigma^2S^2V_{SS}+(r-\delta)SV_S-rV=0$, with solution
$C=Se^{-\delta T}\Phi(d_1)-Ke^{-rT}\Phi(d_2)$, $d_{1,2}=[\ln(S/K)+(r-\delta\pm\tfrac12\sigma^2)T]/\sigma\sqrt T$.
$\Phi(d_2)$ is the risk-neutral exercise probability. Put--call parity holds to $3\times10^{-14}$ on $10^4$ random
contracts. Implied volatility solves $\mathrm{BS}(\sigma)=V^{\text{mkt}}$ by Newton from the Manaster--Koehler start~\cite{mk},
with a bracketing fallback, after checking the no-arbitrage bounds (\code{implied\_vol}).

\textbf{Convergence.} In CRR, $\ln S_N=\ln S_0+\sum_i X_i$ with i.i.d.\ $X_i=\pm\sigma\sqrt{\Delta t}$ whose mean and
variance match risk-neutral GBM to $O(\Delta t)$; by the CLT~\cite{hull} $\ln S_T$ becomes normal and the tree price tends to BS.
The error is $O(1/N)$ but oscillates: terminal nodes sit at $S_0e^{(2j-N)\sigma\sqrt{\Delta t}}$, so the strike's
position $\lambda=\mathrm{frac}\big([\ln(K/S_0)+N\sigma\sqrt{\Delta t}]/2\sigma\sqrt{\Delta t}\big)$ between two nodes
changes with $N$. At the money $\lambda=0$ for even $N$ (strike on a node) and $\tfrac12$ for odd $N$.

\section{Numerical results: convergence, corrections, Greeks}
For the baseline ($S_0=K=100$, $r=5\%$, $\sigma=20\%$, $T=1$; $C_0=10.4506$, $P_0=5.5735$), CRR even $N$
underprice and odd $N$ overprice. A least-squares fit of $\log|e_N|$ on $\log N$ over $N\in[50,1000]$ gives
$|e_N|\approx 1.99/N^{1.000}$ (even $N$) and $1.76/N^{1.000}$ (odd $N$) for CRR and $0.34/N^{1.99}$ for LR
(Fig.~\ref{fig:conv}a). We give no confidence interval: $e_N$ is deterministic, so a regression standard error would
only describe how well a line fits the oscillation. Scaled errors
$N|e_N|$ grow with $\sigma$ and are largest near the money and for long maturities (\code{conv\_error\_heatmap}).
Averaging $N$ and $N{+}1$ cancels the parity term; BS-smoothing of the last step removes the payoff kink;
at the money, Richardson extrapolation $2V_{2N}-V_N$ on even $N$ and LR both reach $O(1/N^2)$ (Fig.~\ref{fig:conv}b).
That success of Richardson is special to $K=S_0$: there $\lambda=0$ for every even $N$, so the error is a clean $C/N$
that extrapolation cancels. Off the money $\lambda$ changes with $N$, the $\lambda(1-\lambda)/N$ term differs between
$N$ and $2N$, and Richardson is no better than plain CRR. Table~\ref{tab:runtime} turns this into cost, counting $N$
only once the error \emph{stays} below the tolerance (not at a lucky zero crossing). Only LR is robust across
strikes: 8--66 steps and at most 0.2\,ms for $10^{-4}$, about the cost of one closed-form BS call in Python
(0.16--0.19\,ms). Plain CRR and JR do not hold $10^{-4}$ within 25{,}000 steps ($|e_N|\approx2/N$ implies about
20{,}000 at the money). Tree Greeks
converge to the BS Greeks, $\Gamma$ oscillating with $N$ like the price (\code{tree\_greeks\_vs\_bs}).

\begin{figure}[t]
\centering
\begin{subfigure}{0.49\linewidth}\includegraphics[width=\linewidth]{conv_loglog_error}\caption{CRR, JR, LR errors}\end{subfigure}
\begin{subfigure}{0.49\linewidth}\includegraphics[width=\linewidth]{conv_corrections}\caption{Corrections}\end{subfigure}
\vspace{-4pt}
\caption{Convergence to Black--Scholes at the baseline contract.}\label{fig:conv}
\vspace{-8pt}
\end{figure}

\begin{table}[H]
\centering\small
\begin{tabular}{lrrrrrr}
\toprule
& \multicolumn{3}{c}{$|e_N|<10^{-3}$} & \multicolumn{3}{c}{$|e_N|<10^{-4}$}\\
\cmidrule(lr){2-4}\cmidrule(lr){5-7}
Method & $K{=}90$ & $K{=}100$ & $K{=}110$ & $K{=}90$ & $K{=}100$ & $K{=}110$\\
\midrule
CRR & 1{,}285 & 2{,}184 & 1{,}964 & $>$25k & $>$25k & $>$25k\\
Jarrow--Rudd & 1{,}285 & 1{,}767 & 1{,}767 & $>$25k & $>$25k & $>$25k\\
CRR avg$(N,N{+}1)$ & 612 & 139 & 551 & 5{,}669 & 1{,}285 & 4{,}586\\
BS-smoothed CRR & 401 & 551 & 236 & 3{,}710 & 5{,}669 & 2{,}428\\
Richardson (even $N$) & 1{,}285 & 11 & 3{,}337 & $>$25k & 39 & $>$25k\\
Leisen--Reimer & 8 & 19 & 19 & 28 & 59 & 66\\
\bottomrule
\end{tabular}
\caption{Steps $N$ after which the call error stays below the tolerance ($S_0=100$, $r=5\%$, $\sigma=20\%$, $T=1$;
grid up to 25{,}000 steps). Runtimes are in \code{conv\_runtime.csv}.}\label{tab:runtime}
\vspace{-10pt}
\end{table}

\section{Real market: smiles, errors, model vs input}
\textbf{Data.} NSE F\&O bhavcopies~\cite{data} for NIFTY, Reliance (RIL) and HUL on 14 days chosen by NIFTY realised volatility:
\emph{calm} (14--18 Jul 2025), \emph{normal} (2--8 Jan 2025), \emph{volatile} (7--11 Apr 2025). The risk-free rate is
the RBI 91-day T-bill cut-off yield before each regime (normal 6.60\%, volatile 6.30\%, calm 5.39\%), as $r=\ln(1+y)$.
\emph{Cleaning} (\code{data\_filter\_summary}). An audit of the raw files changed what counts as a market price. When a
contract does not trade near the close, NSE reports a \emph{theoretical} settlement price from its own model
($\text{settlement}\ne\text{close}$); comparing our models with NSE's model would be circular, so we keep only
settlements set by trades. We then keep the three nearest expiries with $\ge3$ days, volume $\ge10$, open interest
$\ge50$, premium $\ge$ Rs~1 (one tick $\le5\%$), $0.8\le K/F\le1.2$, and prices free of static arbitrage: above the
discounted intrinsic value on the forward, and monotone, convex and with slope $|\partial V/\partial K|\le e^{-rT}$
across strikes. Settlements are closing-half-hour VWAPs, so neighbouring strikes can be struck at different index
levels (worst on 7 Apr 2025). In each violation beyond max(2 ticks, 1\% of price) we drop the lowest-volume
contract. Of 28{,}114 option rows, 6{,}052 remain (72--87\% of the pre-cleaning count in every underlying--regime
cell, all 14 days kept), and implied volatility is solved for all of them. Underlying histories drop two
one-hour Saturday test sessions; close-to-close volatilities use dividend-adjusted closes. Using spot $S$ leaves a
jump between put and call smiles at $K/S=1$; replacing $S$ by the forward implied by parity,
$F=K+e^{rT}(C_K-P_K)$ (median of 5 strikes), and pricing with Black-76~\cite{b76} closes it (\code{mkt\_forward\_fix}).

\textbf{Smiles.} Every regime shows a skewed smile far from the constant $\sigma$ of both models
(Fig.~\ref{fig:smile}). Fitting $\sigma(m)=a+b(m-1)+c(m-1)^2$, $m=K/F$, NIFTY's ATM level $a$ moves 10.5\%
$\to$ 13.5\% $\to$ 21.8\% (calm, normal, volatile) and its skew $b$ steepens from $-0.21$ to $-0.69$ in April.

\begin{figure}[t]
\centering\includegraphics[width=\linewidth]{mkt_smile_NIFTY}
\vspace{-14pt}
\caption{NIFTY implied-volatility smiles (OTM options, Black-76 on the implied forward), first day of each regime;
dashed: 30-day historical volatility.}\label{fig:smile}
\vspace{-6pt}
\end{figure}

\textbf{Numerical error vs input error vs model error.} The tree and BS are not rival models: CRR is a
discretisation of the \emph{same} constant-$\sigma$ GBM, so their gap is numerical error, which falls like $1/N$.
Fig.~\ref{fig:headline} separates three sources of error. Contracts on one date share the spot, forward, volatility
inputs and smile, so they are not independent; every interval below is a 95\% bootstrap CI that resamples whole
trading dates (14 overall, only 4--5 per regime, so the intervals are rough and if anything too narrow).
(i)~\emph{Numerical:} with $\sigma$ fixed, BS and a 200-step CRR tree differ by only 0.02--0.12\,bp of the underlying
(all contracts: 0.05\,bp, CI $[0.04,0.06]$). (ii)~\emph{Input:} switching between five reasonable volatility
estimates (30/60/252-day historical, EWMA, Parkinson) moves prices by 4--21\,bp (all: 11.2, $[9.9,12.5]$). With
$\sigma$ = 30-day historical volatility both methods miss settlement prices by 5--35\,bp (all: 14.9, $[9.8,21.5]$).
The volatile regime is clearly worst for NIFTY (26\,bp $[16,41]$ vs 5--7\,bp in the other regimes); for RIL and HUL
its intervals overlap the calm regime's. (iii)~\emph{Model:} even the best single $\sigma$ (the same-day ATM implied
volatility, so an in-sample fit) leaves 2--15\,bp (all: 6.3, $[4.1,8.8]$). That residual is the smile, a
misspecification that BS and the tree share, because both assume one constant $\sigma$. Across all 6{,}052 contracts
the BS and CRR mean absolute errors agree to 0.01 for every volatility input. \textbf{The numerical gap between the two methods is more than 100 times smaller than the effect of the
volatility input (111--467$\times$ across the nine underlying--regime cells) and 47--237 times smaller than the
constant-$\sigma$ model's own error (the smile).} The ratio depends on $N$: at $N=50$ the numerical gap is about four
times larger.

\begin{figure}[t]
\centering\includegraphics[width=\linewidth]{mkt_model_vs_input}
\vspace{-14pt}
\caption{Mean absolute price differences (bp of the underlying, log scale) per underlying and regime; whiskers are
95\% bootstrap CIs over trading dates.}\label{fig:headline}
\vspace{-8pt}
\end{figure}

\section{Flexibility and model risk}
\textbf{Volatility error vs discretisation error.} At the baseline a 1\% (relative) error in $\sigma$ moves the call
by 0.075, already twice the CRR error at $N=50$ (0.037) and eight times that at $N=200$ (0.009); a 5\% error moves
it by 0.375 (\code{mr\_sigma\_vs\_N}).
\textbf{Time-varying volatility.} Building $\sigma(t)$ from the NIFTY ATM term structure on 7 Apr 2025 and using
Amin's steps~\cite{amin} ($\sigma_i^2\Delta t_i$ constant, so the tree recombines), the tree agrees with BS at the average
variance $\bar\sigma^2=\frac1T\int\sigma^2dt$ up to the same discretisation error as a constant-$\sigma$ tree
(e.g.\ $-0.074$ vs $-0.078$ index points at $N=800$): for Europeans only total variance matters.
\textbf{Discrete dividends.} Escrowed BS~\cite{hull} ($S\mapsto S-\sum_iD_ie^{-rt_i}$) applies $\sigma$ only to the risky part.
In the cash-dividend model the whole price diffuses and drops by $D$ on the ex-date; the tree handles this directly
(CRR to the ex-date, drop $D$ at every node, then BS: \code{binomial\_price\_cash\_div}, checked against
Gauss--Hermite quadrature to $5\times10^{-4}$). For $S=100$, $D=3$ at $t=0.5$, escrowed BS underprices the call by
0.05--0.12 (1.3\% at the money), the case where the tree's flexibility matters (\code{mr\_dividends}). A
\emph{proportional} drop with $\delta=D/\mathbb E[S_{t}]$ is, for Europeans, identical to the escrowed model, and so is
a continuous yield $q=-\ln(1-\mathrm{PV}(D)/S)/T$ with the same forward (BS depends on $S,q$ only through $Se^{-qT}$);
the naive $q=D/S$ removes only about $DT$ and is wrong for short options. On
Reliance options spanning the Rs~5.50 dividend (ex 14 Aug 2025; 159 contracts, but only 6 date--expiry forwards on
5 days), the parity forward with the T-bill carry implies a dividend of $-0.03$ (95\% date-bootstrap CI $[-2.3,1.5]$),
i.e.\ none. The T-bill rate is the wrong carry for a single stock: Reliance expiries with no dividend and $\ge30$
days imply $\ln(F/S)/T$ about 1.3\,pp above it (13 expiries, s.d.\ 1.0\,pp; expiries under 30 days are too noisy,
s.d.\ 3.8\,pp). With that calibrated carry the implied dividend is Rs~2.3 (CI $[0.6,4.3]$). This is closer to 5.50 but
hinges on assuming the spread is the same in July as in January and April. July's own carry, measurable only on a
15-day expiry, would imply Rs~9.7, and a spread of 2.7\,pp would reproduce 5.50 exactly. So the dividend cannot be
identified from parity at this precision, and the pricing errors cannot separate the treatments either. Every
treatment underprices by Rs~0.7--1.1 on average under either carry (MAE 1.0--2.1), for two reasons. One ATM
$\sigma$ across strikes underprices OTM puts (the skew). And with the full Rs~5.50 the model forward sits below the
market's, so calls come out too cheap and ITM puts too dear (\code{mr\_dividends\_market}, \code{mr\_carry\_by\_expiry}).
HUL had no traded options spanning its June ex-date.

\section{Machine learning: can ML choose a better $\sigma$?}
Since the volatility input drives pricing error, we ask whether ML can combine the estimates we already have into a
better $\sigma$ than any single one (\code{src/ml/sigma\_blend.py}). \emph{Inputs} (one row per underlying and day,
all known at the close, no option prices): the 30/60/252-day historical, EWMA and Parkinson~\cite{ewma} volatilities and India
VIX. \emph{Target:} that day's ATM implied volatility of the nearest expiry with $\ge7$ days, i.e.\ the $\sigma$ the
market uses. \emph{Models:} a linear regression (LR) and a random forest (RF)~\cite{rf} on all six inputs, and a linear
regression on 30-day historical and VIX only. Each predicts days it was not trained on, in two ways: \emph{unseen
day} (trained on the other 13 days, so it has seen the same regime; a best case) and \emph{unseen regime}
(trained on the other two regimes). The predicted $\sigma$ prices every contract with Black-76, exactly like the
single estimates.

\begin{figure}[H]
\centering\includegraphics[width=0.92\linewidth]{ml_sigma_blend}
\vspace{-6pt}
\caption{Pricing error for each choice of $\sigma$ (6{,}052 contracts). Whiskers: 95\% CIs over trading dates; they
are wide because all choices share the day-to-day variation, so differences are tested in pairs (text). Dotted line:
best single estimate (EWMA).}\label{fig:ml}
\vspace{-8pt}
\end{figure}

\textbf{Result} (Fig.~\ref{fig:ml}, \code{ml\_sigma\_blend\_pricing}). The best single estimate is EWMA (12.7\,bp); the
usual 30-day historical gives 14.9\,bp. On unseen days the RF blend cuts the error to \textbf{10.1\,bp}: $-4.8$\,bp
against 30-day historical (paired 95\% CI $[-7.4,-2.5]$) and $-2.6$\,bp against EWMA ($[-4.7,-0.7]$), so ML does pick
a better $\sigma$ when it has seen similar market conditions. In a regime it has never seen, both six-input models
fail (19.5 and 21.7\,bp, significantly worse than EWMA): with 42 rows they fit the two training regimes and do not
extrapolate. Only the two-input regression holds (12.2\,bp, level with EWMA; difference $[-4.0,2.1]$); it was added
after seeing the larger models overfit, so treat it as indicative. Same-day ATM implied volatility, which uses option
prices, gives 6.3\,bp: the room left for any $\sigma$ built without option prices. \textbf{ML improves the choice of
$\sigma$ by 21\% (vs EWMA) to 32\% (vs 30-day) within known market conditions, but not in a new regime; there, a simple blend of 30-day
historical volatility and VIX, or EWMA alone, is the safer choice.}

\section{Robustness and limitations}
\textbf{Liquidity filters} (\code{src/market/robustness.py}, \code{rob\_liquidity}). Rebuilding the dataset with
looser filters (volume $\ge1$, no open-interest limit, with or without the Rs~1 minimum premium: 6{,}271 and
6{,}547 contracts) and a stricter one (volume $\ge100$, open interest $\ge500$: 5{,}343) leaves the decomposition
unchanged: numerical 0.047--0.052\,bp, input choice 10.4--11.2\,bp, market error with 30-day historical
14.3--14.9\,bp and with ATM IV 6.1--6.3\,bp, so the input effect stays 215--222 times the numerical gap. The loose
filters add few contracts because trade-set settlements already require trading; that filter, the arbitrage
cleaning and the maturity and moneyness limits stay on in every variant, since without them the ``market price''
is NSE's own model or an arbitrage.
\textbf{Trading-day clock} (\code{rob\_clock}). Volatilities are annualised over 252 trading days while $T$ counts
calendar days. Repricing with variance $\sigma^2 n_{\text{td}}/252$, where $n_{\text{td}}$ counts trading-day returns
to expiry (weekends and NSE holidays excluded, the 1 Feb 2025 Saturday session included), does not help: the error
with 30-day historical volatility moves by $+0.2$\,bp for expiries $\le7$ days (paired CI $[-1.3,2.1]$), $+1.8$ for
8--30 days ($[-0.3,3.9]$) and $+0.3$ beyond ($[-1.7,2.6]$); with EWMA it is $+2.1$ for 8--30 days ($[0.2,4.2]$),
slightly worse, and $+1.8$ overall ($[-0.01,3.7]$). The calendar clock is therefore kept.
Prices with ATM IV do not depend on the clock, since that $\sigma$ is fitted on the same day.

\textbf{Sample.} The market study covers 14 trading days (4--5 per regime) on three underlyings. Date-bootstrap
intervals from so few clusters are rough and probably too narrow, and the ML step learns from only 42 rows; the
direction of every result is consistent across cells, but its size over a full year is untested.
\textbf{Market prices.} Settlements are closing-half-hour VWAPs, not quotes synchronous with the underlying. NSE
holidays after the end of our NIFTY history (1 Aug 2025) are not known to the trading-day count, which affects only
the calm regime's later expiries.
\textbf{Inputs.} The T-bill rate is not the true carry for single stocks (Section~5), and India VIX measures NIFTY
volatility yet is used for Reliance and HUL.
\textbf{Scope.} All options are European, so the tree's main practical advantage (early exercise) is argued, not
measured. Both models use one constant $\sigma$ per day, so the smile residual (6.3\,bp) is outside what either can
fix. The \emph{unseen day} ML test trains on later days of the same regime, so it is a best case; the two-input
regression was chosen after seeing the larger models overfit. Runtimes are single-machine Python timings, so only
their ratios carry over.

\section{Conclusion: which model to choose, and when}
Both models are correct implementations of the same constant-$\sigma$ GBM theory: the tree converges to BS at
$O(1/N)$ with a predictable oscillation, and on 6{,}052 NSE contracts the two differ by 0.05\,bp, against 11.2\,bp for
the choice of $\sigma$ and 6.3\,bp for the smile that neither can produce. So the choice is decided by the
\emph{contract} and by the \emph{$\sigma$ available}, not by which model fits the market better. Table~\ref{tab:rule}
gives the rule that follows from our results.

\begin{table}[H]
\raggedright\small
\begin{tabular}{@{}p{0.27\linewidth}p{0.25\linewidth}p{0.42\linewidth}@{}}
\toprule
Situation & Choose & Evidence from this study\\
\midrule
European option, no cash dividend before expiry & \textbf{Black--Scholes} (Black-76 on the parity forward) & Exact and
0.16--0.19\,ms per price; a tree only adds numerical error (0.05\,bp on NSE data, 0.009 at $N=200$ on the baseline).\\
Volatility varies over time (term structure) & \textbf{Black--Scholes} with the average variance & The time-varying
tree equals BS at $\bar\sigma^2$ to discretisation error ($-0.074$ vs $-0.078$): only total variance matters.\\
Known cash dividend before expiry, priced from spot & \textbf{Tree} with a cash drop at the ex-date & Escrowed BS
underprices the call by up to 1.3\% at the money (synthetic $D=3$); equivalent yield is exact only through the
forward.\\
Dividend already in a market forward & \textbf{Black--Scholes} on that forward & On Reliance (Rs 5.50) all
treatments miss by the same Rs 0.7--1.1: market noise and carry uncertainty swamp the model difference.\\
\raggedright Early exercise or path dependence (American, barrier) & \textbf{Tree} & BS has no closed form; not tested here (NSE
options are European).\\
Any case where a tree is used & \textbf{Leisen--Reimer}, odd $N$ & $10^{-4}$ accuracy in 28--66 steps ($\le0.2$\,ms) at
every strike; CRR/JR need $>25{,}000$; Richardson works only at the money.\\
\midrule
Choosing $\sigma$: option prices exist & Same-day implied volatility & 6.3\,bp, the best constant $\sigma$.\\
\quad no option prices, regime seen before & ML blend (random forest) & 10.1\,bp vs 12.7 (EWMA) and 14.9 (30-day).\\
\quad no option prices, new regime & EWMA or 30-day + VIX regression & 12.2--12.7\,bp; six-input ML overfits (19.5--21.7).\\
\bottomrule
\end{tabular}
\caption{Decision rule. The model rows decide only the last 0.05\,bp; the $\sigma$ rows decide 5--35\,bp.}\label{tab:rule}
\vspace{-8pt}
\end{table}

\textbf{In short: use Black--Scholes for European options on Indian underlyings, with $\sigma$ taken from implied
volatility (or, without option prices, from EWMA or an ML blend); switch to a Leisen--Reimer tree only when the
contract has a cash dividend you must model from spot, early exercise or path dependence.} The error left after the
best constant $\sigma$ is the smile, which both models share; a model with a smile (local or stochastic volatility,
jumps) is the natural next step.

\vspace{-4pt}
{\small
\begin{thebibliography}{99}\setlength{\itemsep}{0pt}
\bibitem{bs} F. Black, M. Scholes, The Pricing of Options and Corporate Liabilities, \emph{J. Political Economy} 81(3), 1973.
\bibitem{crr} J. C. Cox, S. A. Ross, M. Rubinstein, Option Pricing: A Simplified Approach, \emph{J. Financial Economics} 7(3), 1979.
\bibitem{jr} R. A. Jarrow, A. Rudd, \emph{Option Pricing}, Irwin, 1983.
\bibitem{lr} D. P. J. Leisen, M. Reimer, Binomial Models for Option Valuation, \emph{Applied Math. Finance} 3(4), 1996.
\bibitem{amin} K. I. Amin, On the Computation of Continuous Time Option Prices Using Discrete Approximations, \emph{J. Financial Quant. Analysis} 26(4), 1991.
\bibitem{hull} J. C. Hull, \emph{Options, Futures, and Other Derivatives}, 11th ed., Pearson, 2021; S. E. Shreve, \emph{Stochastic Calculus for Finance II}, Springer, 2004.
\bibitem{b76} F. Black, The Pricing of Commodity Contracts, \emph{J. Financial Economics} 3, 1976.
\bibitem{mk} S. Manaster, G. Koehler, The Calculation of Implied Variances from the Black--Scholes Model: A Note, \emph{J. Finance} 37(1), 1982.
\bibitem{ewma} J.P. Morgan, \emph{RiskMetrics Technical Document}, 1996; M. Parkinson, The Extreme Value Method for Estimating the Variance of the Rate of Return, \emph{J. Business} 53(1), 1980.
\bibitem{rf} L. Breiman, Random Forests, \emph{Machine Learning} 45, 2001; F. Pedregosa et al., Scikit-learn, \emph{JMLR} 12, 2011.
\bibitem{data} NSE F\&O bhavcopies, equity price archives and corporate actions (nseindia.com, downloaded 31 Oct and 6 Nov 2026);
NIFTY 50 history (niftyindices.com, 31 Oct 2026); RBI T-bill auction results, press releases 2024-25/1836, 2025-26/15, 2025-26/681 (rbi.org.in).
\end{thebibliography}}
\end{document}
